\documentclass[10pt,a4paper]{amsart}
\usepackage{amsmath,amssymb,amsthm,graphicx}
% Retained preamble for the Stacks Project chapter {CH} (shared project preamble).
% The chapter file pinned-original-source.tex only carries \input{{preamble}} on its first line,
% so no preamble text lives inside the chapter itself.  The Stacks Project uses one single shared
% preamble for every chapter: the snapshot stages it three times (sources/stacks-chow/preamble.tex,
% sources/stacks-etale-cohomology/preamble.tex and focused-seeds/AGP/sources/stacks-snapshot-20261011/
% preamble.tex) and those copies are identical in size.  The block below is copied mechanically,
% line for line, from sources/stacks-chow/preamble.tex lines 14-18 and 154-217: the Xy-pic setup and
% the native theorem declarations together with the mathematical macros.  Nothing was added, removed
% or redefined; the wrapper class and the amsmath/amssymb/amsthm/graphicx base packages are supplied
% by the assembler.
% For commutative diagrams we use Xy-pic
% Chapter: stacks-spaces-more-morphisms.tex -- More on Morphisms of Spaces.
\usepackage[all]{xy}

% We use 2cell for 2-commutative diagrams.
\xyoption{2cell}

% Theorem environments.
%
\theoremstyle{plain}
\newtheorem{theorem}[subsection]{Theorem}
\newtheorem{proposition}[subsection]{Proposition}
\newtheorem{lemma}[subsection]{Lemma}

\theoremstyle{definition}
\newtheorem{definition}[subsection]{Definition}
\newtheorem{example}[subsection]{Example}
\newtheorem{exercise}[subsection]{Exercise}
\newtheorem{situation}[subsection]{Situation}

\theoremstyle{remark}
\newtheorem{remark}[subsection]{Remark}
\newtheorem{remarks}[subsection]{Remarks}

\numberwithin{equation}{subsection}

% Macros
%
\def\lim{\mathop{\mathrm{lim}}\nolimits}
\def\colim{\mathop{\mathrm{colim}}\nolimits}
\def\Spec{\mathop{\mathrm{Spec}}}
\def\Hom{\mathop{\mathrm{Hom}}\nolimits}
\def\Ext{\mathop{\mathrm{Ext}}\nolimits}
\def\SheafHom{\mathop{\mathcal{H}\!\mathit{om}}\nolimits}
\def\SheafExt{\mathop{\mathcal{E}\!\mathit{xt}}\nolimits}
\def\Sch{\mathit{Sch}}
\def\Mor{\mathop{\mathrm{Mor}}\nolimits}
\def\Ob{\mathop{\mathrm{Ob}}\nolimits}
\def\Sh{\mathop{\mathit{Sh}}\nolimits}
\def\NL{\mathop{N\!L}\nolimits}
\def\CH{\mathop{\mathrm{CH}}\nolimits}
\def\proetale{{pro\text{-}\acute{e}tale}}
\def\etale{{\acute{e}tale}}
\def\QCoh{\mathit{QCoh}}
\def\Ker{\mathop{\mathrm{Ker}}}
\def\Im{\mathop{\mathrm{Im}}}
\def\Coker{\mathop{\mathrm{Coker}}}
\def\Coim{\mathop{\mathrm{Coim}}}

% Boxtimes
%
\DeclareMathSymbol{\boxtimes}{\mathbin}{AMSa}{"02}

%
% Macros for moduli stacks/spaces
%
\def\QCohstack{\mathcal{QC}\!\mathit{oh}}
\def\Cohstack{\mathcal{C}\!\mathit{oh}}
\def\Spacesstack{\mathcal{S}\!\mathit{paces}}
\def\Quotfunctor{\mathrm{Quot}}
\def\Hilbfunctor{\mathrm{Hilb}}
\def\Curvesstack{\mathcal{C}\!\mathit{urves}}
\def\Polarizedstack{\mathcal{P}\!\mathit{olarized}}
\def\Complexesstack{\mathcal{C}\!\mathit{omplexes}}
% \Pic is the operator that assigns to X its picard group, usage \Pic(X)
% \Picardstack_{X/B} denotes the Picard stack of X over B
% \Picardfunctor_{X/B} denotes the Picard functor of X over B
\def\Pic{\mathop{\mathrm{Pic}}\nolimits}
\def\Picardstack{\mathcal{P}\!\mathit{ic}}
\def\Picardfunctor{\mathrm{Pic}}
\def\Deformationcategory{\mathcal{D}\!\mathit{ef}}
\usepackage{cleveref}
\crefname{originallemma}{original lemma}{original lemmas}
\Crefname{originallemma}{Original lemma}{Original lemmas}
\makeatletter
\newlabel{lemma-etale-splits-off-quasi-finite-part}{{082H}{source}{}{sourcepointer-1}{}}
\newlabel{lemma-etale-splits-off-quasi-finite-part@cref}{{[originallemma][0][]082H}{[1][0][]source}{}{}{}}
\makeatother
\crefname{originallemma}{original lemma}{original lemmas}
\Crefname{originallemma}{Original lemma}{Original lemmas}
\makeatletter
\newlabel{lemma-normalization-smooth-localization}{{082F}{source}{}{sourcepointer-2}{}}
\newlabel{lemma-normalization-smooth-localization@cref}{{[originallemma][0][]082F}{[1][0][]source}{}{}{}}
\makeatother
\crefname{originalreference}{original reference}{original references}
\Crefname{originalreference}{Original reference}{Original references}
\makeatletter
\newlabel{morphisms-lemma-algebraic-residue-field-extension-closed-point-fibre}{{01TE}{source}{}{sourcepointer-3}{}}
\newlabel{morphisms-lemma-algebraic-residue-field-extension-closed-point-fibre@cref}{{[originalreference][0][]01TE}{[1][0][]source}{}{}{}}
\makeatother
\crefname{originalreference}{original reference}{original references}
\Crefname{originalreference}{Original reference}{Original references}
\makeatletter
\newlabel{morphisms-lemma-locally-quasi-finite-fibres}{{06RT}{source}{}{sourcepointer-4}{}}
\newlabel{morphisms-lemma-locally-quasi-finite-fibres@cref}{{[originalreference][0][]06RT}{[1][0][]source}{}{}{}}
\makeatother
\crefname{originalreference}{original reference}{original references}
\Crefname{originalreference}{Original reference}{Original references}
\makeatletter
\newlabel{spaces-descent-lemma-descending-property-isomorphism}{{041Y}{source}{}{sourcepointer-5}{}}
\newlabel{spaces-descent-lemma-descending-property-isomorphism@cref}{{[originalreference][0][]041Y}{[1][0][]source}{}{}{}}
\makeatother
\crefname{originalreference}{original reference}{original references}
\Crefname{originalreference}{Original reference}{Original references}
\makeatletter
\newlabel{spaces-morphisms-lemma-finite-integral}{{0414}{source}{}{sourcepointer-6}{}}
\newlabel{spaces-morphisms-lemma-finite-integral@cref}{{[originalreference][0][]0414}{[1][0][]source}{}{}{}}
\makeatother
\crefname{originalreference}{original reference}{original references}
\Crefname{originalreference}{Original reference}{Original references}
\makeatletter
\newlabel{spaces-morphisms-lemma-locally-finite-type-quasi-finite-part}{{04NW}{source}{}{sourcepointer-7}{}}
\newlabel{spaces-morphisms-lemma-locally-finite-type-quasi-finite-part@cref}{{[originalreference][0][]04NW}{[1][0][]source}{}{}{}}
\makeatother
\crefname{originalreference}{original reference}{original references}
\Crefname{originalreference}{Original reference}{Original references}
\makeatletter
\newlabel{spaces-morphisms-lemma-normalization-in-disjoint-union}{{0824}{source}{}{sourcepointer-8}{}}
\newlabel{spaces-morphisms-lemma-normalization-in-disjoint-union@cref}{{[originalreference][0][]0824}{[1][0][]source}{}{}{}}
\makeatother
\crefname{originalreference}{original reference}{original references}
\Crefname{originalreference}{Original reference}{Original references}
\makeatletter
\newlabel{spaces-morphisms-lemma-normalization-in-integral}{{0825}{source}{}{sourcepointer-9}{}}
\newlabel{spaces-morphisms-lemma-normalization-in-integral@cref}{{[originalreference][0][]0825}{[1][0][]source}{}{}{}}
\makeatother
\title{Normalization of a finite type separated morphism of algebraic spaces is an isomorphism over an open subspace}
\author{The Stacks Project Authors}
\date{stacks-project:spaces-more-morphisms.tex:local-snapshot-20261011; GNU FDL 1.2 or later}
\begin{document}
\maketitle
\noindent\textit{Editorial note.} Counted span is the author's own proof environment, source lines 7611-7707, immediately following the statement at 7592-7609.\par\medskip
\noindent\textit{Editorial note.} Seven cross-chapter pointers (Morphisms of Spaces, Descent on Spaces) plus two same-chapter lemmas are recorded in assembly.external\_references.\par\medskip
\noindent\textit{Editorial note.} Same-name pair note: this is the more general (algebraic spaces) side; the schemes side is a separate statement in the other chapter and is not counted with this card.\par\medskip
\noindent\textit{Editorial note.} Steps the author leaves to the reader are flagged here for the reviewer: source line(s) 7632 carry an "omitted"/"follows easily"/"immediate from the definitions" phrase. Those sentences are retained verbatim; the displayed argument around them stays inside the counted span.\par\medskip
\noindent\textit{Allowed standard prerequisites (editorial).}\begin{itemize}
\item Etale splitting off of the quasi-finite part. chapter lemma lemma-etale-splits-off-quasi-finite-part, cited at proof line 7694. Produces, etale locally over a point, a decomposition into a finite part and a part whose normalization is followed.
\item Localization of the smooth/normalization comparison. chapter lemma lemma-normalization-smooth-localization, cited in the proof. Lets the normalization be computed after localizing, which is what the local step needs.
\item Normalization commutes with disjoint unions and with integral morphisms. Morphisms of Spaces, lemmas spaces-morphisms-lemma-normalization-in-disjoint-union and spaces-morphisms-lemma-normalization-in-integral, cited at proof lines 7697-7701 onward. Identifies the normalization of Y in X = V coproduct W with V coproduct W', which yields the open subspace U'.
\item Residue field extension realised by a closed point of a fibre. Morphisms of Spaces, lemmas spaces-morphisms-lemma-algebraic-residue-field-extension-closed-point-fibre and spaces-morphisms-lemma-locally-quasi-finite-fibres, cited in the proof. Supplies the point over which the etale-local splitting is performed.
\end{itemize}\medskip
\ifdefined\hypertarget\hypertarget{sourcepointer-1}{}\fi
\noindent\textit{Source reference.} Original lemma 082H: sources/stacks-spaces-more-morphisms/pinned-original-source.tex, source lines 7472--7472. Navigation pointer only: the counted proof cites this statement by \textbackslash{}ref as a standing input; the deduction the counted proof has to make is carried out inside the counted span, and no part of that deduction is delegated to this pointer.\par\medskip
\ifdefined\hypertarget\hypertarget{sourcepointer-2}{}\fi
\noindent\textit{Source reference.} Original lemma 082F: sources/stacks-spaces-more-morphisms/pinned-original-source.tex, source lines 6141--6141. Navigation pointer only: the counted proof cites this statement by \textbackslash{}ref as a standing input; the deduction the counted proof has to make is carried out inside the counted span, and no part of that deduction is delegated to this pointer.\par\medskip
\ifdefined\hypertarget\hypertarget{sourcepointer-3}{}\fi
\noindent\textit{Source reference.} Original reference 01TE: focused-seeds/geometry/S9-spaces-finite-type-separated/ref-tags.txt, source lines 4--4. Navigation pointer only: the counted proof cites this statement by \textbackslash{}ref as a standing input; the deduction the counted proof has to make is carried out inside the counted span, and no part of that deduction is delegated to this pointer.\par\medskip
\ifdefined\hypertarget\hypertarget{sourcepointer-4}{}\fi
\noindent\textit{Source reference.} Original reference 06RT: focused-seeds/geometry/S9-spaces-finite-type-separated/ref-tags.txt, source lines 5--5. Navigation pointer only: the counted proof cites this statement by \textbackslash{}ref as a standing input; the deduction the counted proof has to make is carried out inside the counted span, and no part of that deduction is delegated to this pointer.\par\medskip
\ifdefined\hypertarget\hypertarget{sourcepointer-5}{}\fi
\noindent\textit{Source reference.} Original reference 041Y: focused-seeds/geometry/S9-spaces-finite-type-separated/ref-tags.txt, source lines 6--6. Navigation pointer only: the counted proof cites this statement by \textbackslash{}ref as a standing input; the deduction the counted proof has to make is carried out inside the counted span, and no part of that deduction is delegated to this pointer.\par\medskip
\ifdefined\hypertarget\hypertarget{sourcepointer-6}{}\fi
\noindent\textit{Source reference.} Original reference 0414: focused-seeds/geometry/S9-spaces-finite-type-separated/ref-tags.txt, source lines 7--7. Navigation pointer only: the counted proof cites this statement by \textbackslash{}ref as a standing input; the deduction the counted proof has to make is carried out inside the counted span, and no part of that deduction is delegated to this pointer.\par\medskip
\ifdefined\hypertarget\hypertarget{sourcepointer-7}{}\fi
\noindent\textit{Source reference.} Original reference 04NW: focused-seeds/geometry/S9-spaces-finite-type-separated/ref-tags.txt, source lines 8--8. Navigation pointer only: the counted proof cites this statement by \textbackslash{}ref as a standing input; the deduction the counted proof has to make is carried out inside the counted span, and no part of that deduction is delegated to this pointer.\par\medskip
\ifdefined\hypertarget\hypertarget{sourcepointer-8}{}\fi
\noindent\textit{Source reference.} Original reference 0824: focused-seeds/geometry/S9-spaces-finite-type-separated/ref-tags.txt, source lines 9--9. Navigation pointer only: the counted proof cites this statement by \textbackslash{}ref as a standing input; the deduction the counted proof has to make is carried out inside the counted span, and no part of that deduction is delegated to this pointer.\par\medskip
\ifdefined\hypertarget\hypertarget{sourcepointer-9}{}\fi
\noindent\textit{Source reference.} Original reference 0825: focused-seeds/geometry/S9-spaces-finite-type-separated/ref-tags.txt, source lines 10--10. Navigation pointer only: the counted proof cites this statement by \textbackslash{}ref as a standing input; the deduction the counted proof has to make is carried out inside the counted span, and no part of that deduction is delegated to this pointer.\par\medskip
\begin{lemma}
\label{lemma-finite-type-separated}
Let $S$ be a scheme. Let $f : X \to Y$ be a morphism of algebraic
spaces over $S$ which is of finite type and separated.
Let $Y'$ be the normalization of $Y$ in $X$. Picture:
$$
\xymatrix{
X \ar[rd]_f \ar[rr]_{f'} & & Y' \ar[ld]^\nu \\
& Y &
}
$$
Then there exists an open subspace $U' \subset Y'$ such that
\begin{enumerate}
\item $(f')^{-1}(U') \to U'$ is an isomorphism, and
\item $(f')^{-1}(U') \subset X$ is the set of points at which
$f$ is quasi-finite.
\end{enumerate}
\end{lemma}

\begin{proof}
By Morphisms of Spaces, Lemma
\ref{spaces-morphisms-lemma-locally-finite-type-quasi-finite-part}
there is an open subspace $U \subset X$ corresponding to the points
of $|X|$ where $f$ is quasi-finite. We have to prove
\begin{enumerate}
\item[(a)] the image of $|U| \to |Y'|$ is $|U'|$ for some open subspace
$U'$ of $Y'$,
\item[(b)] $U = f^{-1}(U')$, and
\item[(c)] $U \to U'$ is an isomorphism.
\end{enumerate}
Since formation of $U$ commutes with arbitrary base change
(Morphisms of Spaces, Lemma
\ref{spaces-morphisms-lemma-locally-finite-type-quasi-finite-part}),
since formation of the normalization $Y'$ commutes with smooth
base change (Lemma \ref{lemma-normalization-smooth-localization}),
since \'etale morphisms are open, and
since ``being an isomorphism'' is fpqc local on the base
(Descent on Spaces, Lemma
\ref{spaces-descent-lemma-descending-property-isomorphism}),
it suffices to prove (a), (b), (c) \'etale locally on $Y$
(some details omitted). Thus
we may assume $Y$ is an affine scheme. This implies that $Y'$ is
an (affine) scheme as well.

\medskip\noindent
Let $x \in |U|$. Claim: there exists an open
neighbourhood $f'(x) \in V \subset Y'$ such that $(f')^{-1}V \to V$ is an
isomorphism. We first prove the claim implies the lemma.
Namely, then $(f')^{-1}V \cong V$ is a scheme (as an open
of $Y'$), locally of finite type over $Y$ (as an open subspace of $X$),
and for $v \in V$ the residue field extension
$\kappa(v)/\kappa(\nu(v))$ is algebraic (as
$V \subset Y'$ and $Y'$ is integral over $Y$). Hence the fibres
of $V \to Y$ are discrete (Morphisms, Lemma
\ref{morphisms-lemma-algebraic-residue-field-extension-closed-point-fibre})
and $(f')^{-1}V \to Y$ is locally quasi-finite
(Morphisms, Lemma \ref{morphisms-lemma-locally-quasi-finite-fibres}).
This implies $(f')^{-1}V \subset U$ and $V \subset U'$. Since $x$ was
arbitrary we see that (a), (b), and (c) are true.

\medskip\noindent
Let $y = f(x) \in |Y|$. Let $(T, t) \to (Y, y)$ be an \'etale morphism
of pointed schemes. Denote by a subscript ${}_T$ the base change to $T$.
Let $z \in X_T$ be a point in the fibre $X_t$ lying over $x$.
Note that $U_T \subset X_T$ is the set of points where $f_T$ is
quasi-finite, see Morphisms of Spaces, Lemma
\ref{spaces-morphisms-lemma-locally-finite-type-quasi-finite-part}.
Note that
$$
X_T \xrightarrow{f'_T} Y'_T \xrightarrow{\nu_T} T
$$
is the normalization of $T$ in $X_T$, see
Lemma \ref{lemma-normalization-smooth-localization}.
Suppose that the claim holds for $z \in U_T \subset X_T \to Y'_T \to T$, i.e.,
suppose that we can find an open neighbourhood
$f'_T(z) \in V' \subset Y'_T$ such that $(f'_T)^{-1}V' \to V'$ is an
isomorphism. The morphism $Y'_T \to Y'$ is \'etale hence the image
$V \subset Y'$ of $V'$ is open. Observe that $f'(x) \in V$ as $f'_T(z) \in V'$.
Observe that
$$
\xymatrix{
(f'_T)^{-1}V' \ar[r] \ar[d] & (f')^{-1}(V) \ar[d] \\
V' \ar[r] & V
}
$$
is a fibre square (as $Y'_T \times_{Y'} X = X_T$).
Since the left vertical arrow is an isomorphism
and $\{V' \to V\}$ is a \'etale covering, we conclude that the right vertical
arrow is an isomorphism by
Descent on Spaces, Lemma
\ref{spaces-descent-lemma-descending-property-isomorphism}.
In other words, the claim holds for $x \in U \subset X \to Y' \to Y$.

\medskip\noindent
By the result of the previous paragraph to prove the claim for
$x \in |U|$, we may replace $Y$ by an \'etale neighbourhood $T$ of
$y = f(x)$ and $x$ by any point lying over $x$ in $T \times_Y X$.
Thus we may assume there is a decomposition
$$
X = V \amalg W
$$
into open and closed subspaces where $V \to Y$ is finite and $x \in V$,
see Lemma \ref{lemma-etale-splits-off-quasi-finite-part}.
Since $X$ is a disjoint union of $V$ and $W$ over $Y$ and since
$V \to Y$ is finite we see that the
normalization of $Y$ in $X$ is the morphism
$$
X = V \amalg W \longrightarrow V \amalg W' \longrightarrow S
$$
where $W'$ is the normalization of $Y$ in $W$, see
Morphisms of Spaces, Lemmas
\ref{spaces-morphisms-lemma-normalization-in-disjoint-union},
\ref{spaces-morphisms-lemma-finite-integral}, and
\ref{spaces-morphisms-lemma-normalization-in-integral}.
The claim follows and we win.
\end{proof}

\end{document}
